\chapter{Probabilistic bijective combinatorics for Macdonald polynomials} \label{chap:probabilistic}

This chapter draws upon the research presented in:
\begin{quote}
    \fullcite{DM25}.
\end{quote}

\section{Probabilistic bijections}

Probabilistic bijections,
also referred to as \emph{bijectivization} or \emph{coupling} by \cite{BP19},
generalize weight-preserving bijections.
We follow the exposition given in~\cite{AF22}.

\begin{definition}[Probability map]
	Let \(\algebra\) be an algebra.
	Let \(\mathbf{T}\) and \(\mathbf{U}\) be sets.
	A \vocab{probability map} from \(\mathbf{T}\) to \(\mathbf{U}\) is
	a function \(\prob{} \colon \mathbf{T} \times \mathbf{U} \to \algebra\) such that
	\begin{equation}
		\sum_{U \in\mathbf{U}} \prob{}(T, U) = 1
	\end{equation}
	for every \(T \in \mathbf{T}\).
\end{definition}

Note that these do not formally correspond to probability maps
in the sense of probability theory
because the values of \(\prob{}\)
are not required to be in the interval \([0, 1]\) or even real numbers.

\begin{definition}[Probabilistic bijection {\cite{AF22}}]\label{def:prob}
	Let \(\algebra\) be an algebra.
	Let \(\mathbf{T}\) and \(\mathbf{U}\) be finite sets
	equipped with weight functions
	\(\wt_{\mathbf{T}} \colon \mathbf{T} \to \algebra\) and
	\(\wt_{\mathbf{U}} \colon \mathbf{U} \to \algebra\).
	A \vocab{probabilistic bijection} between
	\( (\mathbf{T}, \wt_{\mathbf{T}}) \) and
	\( (\mathbf{U}, \wt_{\mathbf{U}}) \) is
	a pair of probability maps
	\(\prob{\mathbf{T}} \colon \mathbf{T} \times \mathbf{U} \to \algebra\) and
	\(\prob{\mathbf{U}} \colon \mathbf{U} \times \mathbf{T} \to \algebra\)
	such that
	\begin{equation}
		\wt_{\mathbf{T}}(T) \prob{\mathbf{T}}(T, U)
		= \wt_{\mathbf{U}}(U) \prob{\mathbf{U}}(U, T)
	\end{equation}
	for every \(T \in \mathbf{T}\) and \(U \in \mathbf{U}\).
\end{definition}

A weight-preserving bijection \( f \colon \mathbf{T} \to \mathbf{U} \)
canonically induces a probabilistic bijection
between \((\mathbf{T}, \wt_{\mathbf{T}})\) and \((\mathbf{U}, \wt_{\mathbf{U}})\)
by setting
\begin{equation}
	\prob{\mathbf{T}}(T, U) = \prob{\mathbf{U}}(U, T) =
	\begin{cases}
		1 & \text{if } f(T) = U, \\
		0 & \text{otherwise},
	\end{cases}
\end{equation}
explaining the choice of terminology.
Moreover, as shown in \cref{proposition:equal-weight-sum},
this generalization retains a key property of weight-preserving bijections:
the sum of the weights of \(\mathbf{T}\) and of \(\mathbf{U}\) are equal.

\begin{proposition}
	\label{proposition:equal-weight-sum}
	Let \(\prob{\mathbf{T}} \colon \mathbf{T} \times \mathbf{U} \to \algebra\)
	and \(\prob{\mathbf{U}} \colon \mathbf{U} \times \mathbf{T} \to \algebra\)
	form a probabilistic bijection between
	\( (\mathbf{T}, \wt_{\mathbf{T}}) \) and
	\( (\mathbf{U}, \wt_{\mathbf{U}}) \).
	Then,
	\begin{equation}
		\sum_{T \in \mathbf{T}} \wt(T) = \sum_{U \in\mathbf{U}} \wt(U).
	\end{equation}
\end{proposition}

\begin{proof}
	We compute
	\begin{align}
		\sum_{T \in \mathbf{T}} \wt_{\mathbf{T}}(T)
		&= \sum_{T \in \mathbf{T}} \sum_{U \in \mathbf{U}}
			\wt_{\mathbf{T}}(T) \prob{\mathbf{T}}(T, U) \\
		&= \sum_{U \in \mathbf{U}} \sum_{T \in \mathbf{T}}
			\wt_{\mathbf{U}}(U) \prob{\mathbf{U}}(U, T)
		= \sum_{U \in \mathbf{U}} \wt_{\mathbf{U}}(U). \qedhere
	\end{align}
\end{proof}

A Markov-theoretical interpretation of a probabilistic bijection is that
it defines a bipartite Markov chain on the state space \( \mathbf{T} \cup \mathbf{U} \)
with transition probabilities given by the function \( \prob{} \) under a specific balance condition.
In the case of \cref{def:prob}, this condition is detailed balance,
but a probabilistic bijection can also be defined using a more general balance condition.

\todo[caption={Add the remainder of the stuff from \cite{DM25}.}]{Add the remainder of the stuff from \cite{DM25}. This is mostly copy-paste, but I need to make sure the notation is consistent with the rest of the thesis, and I need to delete definitions that have already been given.}